\documentclass{article}
\usepackage[a4paper, total={6.5in, 10in}]{geometry}
\setlength{\parindent}{0pt}
\usepackage{tcolorbox}
\tcbset{colback=white!94!black, colframe=black!60!white, coltitle = white, boxrule=0.8pt, arc=2mm}
\newtcolorbox{boxs}[2][]{title= #2,#1}
\usepackage{datetime2}
\usepackage[british]{babel}
\usepackage{amsfonts}
\usepackage{amsmath}

\title{\textbf{Advanced Calculus HW 5}}
\author{Robert O'Connell}
\date{\today}
\begin{document}

\vspace{10mm}

\subsection*{Question 2}
\subsubsection*{(a)}
\[
f(x,y,z) = e^{x^2 + y^2} \sin(y^2 + xz^2) + z
\]

\[
\frac{\partial f}{\partial x} = 2x e^{x^2 + y^2} \sin(y^2 + z^2 x) + e^{x^2 + y^2} z^2 \cos(y^2 + z^2 x)
\]
\[
\frac{\partial f}{\partial y} = 2y e^{x^2 + y^2} \sin(y^2 + z^2 x) + e^{x^2 + y^2} 2y \cos(y^2 + z^2x)
\]
\[
\frac{\partial f}{\partial z} = e^{x^2 + y^2} 2z \cos(y^2 + z^2 x) + 1
\]

\subsubsection*{(b)}
\[
f(\rho, \phi, \theta) = \rho^2 cos \theta \sin^2 \phi
\]

\[
\frac{\partial f }{\partial \rho} = 2 \rho \cos \theta \sin^2 \phi 
\]
\[
\frac{\partial f}{\partial \phi} = \rho^2 \cos \theta 2 \cos \phi \sin \phi
\]
\[
\frac{\partial f}{\partial \theta} =  - \rho^2 \sin \theta \sin^2 \phi
\]

\subsubsection*{(c)}
\[
f(x,y,z) = x \ln(x^2z^2 + g^2(y^2z^3))
\]

\[
\frac{\partial f}{\partial x} = \ln(x^2z^2 + g^2(y^2z^3)) + 2xz^2 \frac{x}{x^2z^2 + g^2(y^2z^3)}
\]
\[
\frac{\partial f}{\partial y} = x \frac{4z^3y g'(y^2z^3) g(y^2z^3)}{x^2z^2 + g^2(y^2z^3)}
\]
\[
\frac{\partial f}{\partial z} = x \frac{2zx^2 + (6 y^2 z^2 g'(y^2z^3) g(y^2z^3))}{x^2z^2 + g^2(y^2z^3)}
\]

\subsection*{Question 3}
\[
\lim_{(x,y,x)\to(0,0,0)}{\frac{xyz}{x^2 + y^4 z^4}}
\]

\begin{itemize}
    \item Along the curve \(x = t, y = t, z = t\)
    \[
    \lim_{t\to 0}{\frac{t^3}{t^2 + 2t^4}} = 0
    \]
    \item Along the curve \(x = t, y = \sqrt{t}, z = \sqrt{t}\)
    \[
    \lim_{t\to 0}{\frac{t^2}{3t^2}} = \frac{1}{3}
    \]
\end{itemize}
Since these limits are not equal to each other, we can conlcude that \(f\) is not 
continuous at \((0,0,0)\)

\subsection*{Question 4}
\[
f(x,y) = \cos(\pi x^3 y^4)
\]

\subsubsection*{(a)}
\[
f_x = -3\pi x^2 y^4 \sin(\pi x^3y^4)
\]
\[
f_y= -4\pi y^3x^3 \sin(\pi x^3 y^4)
\]

\subsubsection*{(b)}
This is simply the curve with the \(x\) value fixed at \(x = 2\)
\[
z = \cos(8\pi y^4)
\]

\subsubsection*{(c)}
To find the tangent line to the curve at the point \((2,1,1)\), we need the slope of the line at this point,
which we can use the partial derivative \(f_y\)
\[
f_y (2,1) = 0
\]
Thus the tangent line is simply 
\[
\vec{r}(t) = (2,1,1) + (0,1,0) t
\]



\end{document}